\documentclass[10pt,a4paper]{article}
\usepackage[utf8]{inputenc}
\usepackage[spanish]{babel}
\usepackage{url}
\usepackage{hyperref}
\usepackage{amsmath}
\usepackage{amsfonts}
\usepackage{amssymb}
\usepackage{graphicx}
\usepackage{float}
\usepackage{lipsum}
\usepackage{multicol}
\usepackage{xcolor}
\newcommand{\celda}[1]{
	\begin{minipage}{2.5cm}
		\vspace{5mm}
		#1
		\vspace{5mm}
	\end{minipage}
}
\definecolor{azul}{rgb}{0.0, 0.53, 0.74}
\usepackage[left=2.00cm, right=2.00cm, top=2.00cm, bottom=2.00cm]{geometry}
\author{Aprendiendo \LaTeX\,}
\title{Mi primer paper}

\begin{document}
	
	\begin{figure}[H]
		\raggedright
		\includegraphics[scale=0.1]{sm.png} \hfill \includegraphics[scale=0.35]{villa.png}
	\end{figure}

	\vspace{7.5mm}
	
	\begin{center}
		{\Large \textbf{Mi primera investigación aplicada a la inteligenia artificial}}\\
		\vspace{2mm}
		{\large Manuel Merino$^{1}$, Aprendiendo \LaTeX$^{1}$, Pablo Pérez$^{2}$}\\
		\vspace{7.5mm}
		$^1$\textit{Universidad Nacional Mayor de San Marcos. Facultad de Ciencias Matemáticas.} \\
		$^2$\textit{Universidad Nacional Federico Villareal. Facultad de Ciencias Naturales.}
	\end{center}

	\begin{center}
		\textcolor{azul}{\rule{150mm}{0.5mm}}
	\end{center}		

	\begin{abstract}
		\lipsum[4] \lipsum[5]
		\vspace{2mm}
		\underline{\textbf{Palabras clave:}} \hspace{1mm} \textit{palabra1, palabra2, palabra3.}
	\end{abstract}

	\begin{center}
		\textcolor{azul}{\rule{150mm}{0.5mm}}
	\end{center}		

	\vspace{5mm}
	
	\begin{multicols}{2}
		\section{Introducción}
			\lipsum[2] \lipsum[1] \lipsum[4] \lipsum[5]
			
		\section{Conocimientos Previos}
			\lipsum[3] \lipsum[7]
	\end{multicols}
	
	\begin{table}[H]
		\centering
		\begin{tabular}{cccc}
			\hline
			\celda{Nombres} & \celda{Apellidos} & \celda{Edad} & \celda{País}\\
			\hline
			\celda{Manuel} & \celda{Merino} & \celda{25} & \celda{Perú}\\
			\hline
			\celda{Pablo} & \celda{Perez} & \celda{27} & \celda{México}\\
			\hline
		\end{tabular}
		\caption{Esta es mi primera tabla en mi papaer}
		\label{tb: Tabla1}
	\end{table}

	\begin{multicols}{2}
		\lipsum[1]
		\section{Conociendo más de la Inteligencia Artificial.}
			\lipsum
			
			\begin{figure}[H]
				\centering
				\includegraphics[scale=0.1]{sm.png}
				\caption{Esta es mi primera figura del paper}
				\label{fig: Figura1}
			\end{figure}
		
		\lipsum[2]
		
		\section{Aplicación y Alcances}
			\lipsum
			
	    \begin{thebibliography}{00}
	    
	    \bibitem{Referencia1}
	        \newblock \textit{Título de la referencia 1}
	        \newblock Autor.
	        \newblock Año y editor.
	    
	    \bibitem{Referencia2}
	        \newblock \textit{Título de la referencia 2}
	        \newblock Autor.
	        \newblock Año y editor.
	        
	    \bibitem{Referencia3}
	        \newblock \textit{Título de la referencia 3}
	        \newblock Autor.
	        \newblock Año y editor.
	        
	    \bibitem{Referencia4}
	        \newblock \textit{Título de la referencia 4}
	        \newblock Autor.
	        \newblock Año y editor.
	    
	    \bibitem{Referencia5}
	        \newblock \textit{Título de la referencia 5}
	        \newblock Autor.
	        \newblock Año y editor.
	        
	    \bibitem{Referencia6}
	        \newblock \textit{Título de la referencia 6}
	        \newblock Autor.
	        \newblock Año y editor.
	    
	    \end{thebibliography}
	\end{multicols}
	
	\vspace{3cm}
	Para un mejor entendimiento de la plantilla ver el video en youtube:\\
	
	\url{https://www.youtube.com/watch?v=34xYuLxnzHo}
	
	
\end{document}